\section{Reduction to exponential decay}\label{sec:reduction}

The first step replaces a general $\ell^1$ idempotent by an idempotent
whose coefficients have exponential decay in a finitely generated
subgroup. The replacement preserves every trace coefficient exactly.
This permits the polynomially weighted sums that arise from differential
forms in arbitrarily large, but fixed, dimensions.

For a complex matrix $A=(A_{ij})$, write $|A|_1=\sum_{i,j}|A_{ij}|$.
On $M_q(\ell^1(G))$ we use
\[
 \|a\|_1=\sum_{h\in G}|a(h)|_1.
\]
This norm is submultiplicative. The identity matrix has norm $q$;
none of the power-series arguments below requires its norm to be one.

\begin{lemma}[Traciality]\label{red:trace}
For every conjugacy class $C$ of a discrete group $G$, the functional
\[
 T_C(a)=\sum_{h\in C}\tr a(h),\qquad a\in M_q(\ell^1(G)),
\]
is continuous and satisfies $T_C(ab)=T_C(ba)$. It is therefore invariant
under similarity by invertible matrices over $\ell^1(G)$.
\end{lemma}
\begin{proof}
The bound $|T_C(a)|\le\|a\|_1$ gives continuity. Expand $T_C(ab)$
as the absolutely convergent sum
\[
 \sum_{\substack{x,y\in G\\xy\in C}}\tr\bigl(a(x)b(y)\bigr).
\]
The conditions $xy\in C$ and $yx\in C$ are equivalent, and the scalar
matrix trace is cyclic. Interchanging $x$ and $y$ proves the identity.
If $u$ is invertible, apply it to $ua$ and $u^{-1}$.
\end{proof}

The same rearrangement applies to rectangular matrices: if
$a\in M_{q,r}(\ell^1(G))$ and $b\in M_{r,q}(\ell^1(G))$, then
$T_C(ab)=T_C(ba)$, with the trace in the appropriate size on each side.
An isomorphism between two finitely generated projective modules,
together with its inverse, can be extended by their defining
idempotents to such matrices $a,b$ with products equal to the two
idempotents. This proves the module invariance used in the definition
of $\operatorname{HS}^1$.

\begin{proposition}[Exponentially decaying representative]\label{red:reduction}
Let $e\in M_q(\ell^1(G))$ satisfy $e^2=e$. There are a finitely generated
subgroup $\Gamma\le G$, a word length $|\cdot|$ on $\Gamma$, a number
$c>0$, and an idempotent $p\in M_q(\ell^1(\Gamma))$ such that
\begin{equation}\label{red:exponential}
 \sum_{h\in\Gamma}|p(h)|_1e^{c|h|}<\infty,
\end{equation}
and $p$ is similar to $e$ in $M_q(\ell^1(G))$.
If the trace of every such $p$ has finite support on finite-order
conjugacy classes of $\Gamma$, the same conclusion holds for every $e$
over every discrete group.
\end{proposition}
\begin{proof}
Choose a finite-support matrix $a$ close to $e$ and put
$b=4(a-a^2)$. We may require $\|b\|_1<1$. Let
\[
 Q=(1-b)^{-1/2}=\sum_{j=0}^\infty
       \frac{\binom{2j}{j}}{4^j}b^j,
 \qquad
 p=\frac12 I+\left(a-\frac12 I\right)Q.
\]
The series converges absolutely in the matrix Banach algebra.
All its factors commute with $a$, and the scalar power-series identity
gives $Q^2(1-b)=I$. Since
$(a-I/2)^2=(I-b)/4$, we obtain $(p-I/2)^2=I/4$, hence $p^2=p$.
Moreover $p\to e$ as $a\to e$.

Let $\Gamma$ be generated by the finite union of the supports of the
entries of $a$. Fix a finite symmetric generating set and its word
length. The norm
\[
 \|v\|_c=\sum_{h\in\Gamma}|v(h)|_1e^{c|h|}
\]
is submultiplicative by $|xy|\le |x|+|y|$. The element $b$ has finite
support, so $\|b\|_c\to\|b\|_1$ as $c\downarrow0$.
Choose $c>0$ with $\|b\|_c<1$. The same series defining $Q$ then
converges in this weighted algebra, proving \eqref{red:exponential}.

Finally, set
\[
 u=pe+(1-p)(1-e).
\]
Direct multiplication gives $ue=pu$ and
$u-I=(p-e)(2e-I)$. For $a$ sufficiently close to $e$, the latter
has norm less than one, so $u$ is invertible and $p=ueu^{-1}$.
Lemma~\ref{red:trace} preserves all the trace coefficients.

For a conjugacy class $C$ of $G$, its intersection with $\Gamma$
is a union of conjugacy classes of $\Gamma$. Absolute convergence gives
\[
 T_C(p)=\sum_{\substack{D\in\mathcal C(\Gamma)\\D\subset C}}T_D(p).
\]
Thus passing from $\Gamma$ to $G$ merges classes and cannot enlarge
a finite set of surviving classes. The order of an element does not
change under subgroup inclusion. This proves the final assertion.
\end{proof}

Henceforth $G$ is finitely generated, $d(x,y)=|x^{-1}y|$ is a fixed
left-invariant word metric, and $p$ satisfies \eqref{red:exponential}.
We fix a constant $\kappa_G$ such that
\begin{equation}\label{red:balls}
 |B(x,s)|\le e^{\kappa_G s}\qquad(x\in G,\ s\ge1).
\end{equation}
Such a bound follows by counting words in the generators. It imposes
no additional growth hypothesis on $G$.
Constants without a subscript may depend on this metric and on fixed
construction choices. A constant denoted $C_R$ may also depend on the
radius parameter $R$ introduced later, but will not depend on a tuple,
its twist $g$, or the dimension $n$.
